\chapter{raymath}

raymath ships with raylib but is not included by \file{raylib.h}. You ask for it
separately, and it has its own cheatsheet:

\begin{lstlisting}
#include "raymath.h"
\end{lstlisting}

It is header-only and every function is small and inlined, so using it costs
nothing at run time. Chapter 11 listed the handful you meet first; this is the
whole shape of it.

\section{The three ideas worth understanding}

Everything else is arithmetic. These three are concepts.

\subsection{Normalize: direction without length}

\begin{lstlisting}
Vector3 dir = Vector3Normalize(Vector3Subtract(target, me));
\end{lstlisting}

A normalized vector has length exactly 1, so it carries a \emph{direction} and
nothing else. Multiply it by a speed and you get a velocity; by a distance and
you get an offset. Almost every movement line you will ever write is:

\begin{lstlisting}
position = Vector3Add(position, Vector3Scale(dir, speed*dt));
\end{lstlisting}

\gotcha{Normalizing a zero vector divides by zero. raylib's version guards
against it and returns zero, but the habit of checking
\lstinline|Vector3Length(v) > 0| before normalizing will save you elsewhere ---
and it is why the movement code in chapter 12 tests the length first.}

\subsection{Dot product: how much two directions agree}

For two \emph{normalized} vectors, the dot product is a single number:

\begin{center}
\begin{tabular}{@{}ll@{}}
\toprule
$+1$ & pointing the same way \\
$0$ & perpendicular \\
$-1$ & pointing opposite ways \\
\bottomrule
\end{tabular}
\end{center}

That makes it the answer to a surprising number of game questions:

\begin{lstlisting}
// Is the enemy in front of me, or behind?
Vector3 toEnemy = Vector3Normalize(Vector3Subtract(enemy, camera.position));
float facing = Vector3DotProduct(forward, toEnemy);
if (facing > 0.0f) { /* somewhere in front */ }

// Is it inside my 60-degree cone of vision?
if (facing > cosf(30.0f*DEG2RAD)) { /* it can see me */ }
\end{lstlisting}

That second one is a whole enemy vision system in one line, and it is far cheaper
than computing angles.

\subsection{Cross product: perpendicular to both}

3D only. Give it two directions, get a third at right angles to both:

\begin{lstlisting}
Vector3 right = Vector3CrossProduct(forward, (Vector3){ 0, 1, 0 });
\end{lstlisting}

That is the line from chapter 12 that gives you strafing. It is also how you
compute a surface normal from a triangle's edges, which is how lighting knows
which way a face points.

\gotcha{The cross product is not commutative: \lstinline|cross(a, b)| points the
opposite way to \lstinline|cross(b, a)|. If your strafe keys are swapped, this is
why, and swapping the arguments is the fix.}

\section{The families}

\subsection{Utils}

\begin{center}
\begin{tabular}{@{}ll@{}}
\toprule
\lstinline|Lerp(a, b, t)| & blend: $t=0$ gives a, $t=1$ gives b \\
\lstinline|Clamp(value, min, max)| & keep it in range \\
\lstinline|Normalize(value, start, end)| & map a range onto 0..1 \\
\lstinline|Remap(v, inA, inB, outA, outB)| & map one range onto another \\
\lstinline|Wrap(value, min, max)| & wrap around, for angles \\
\lstinline|FloatEquals(a, b)| & float comparison that is not a lie \\
\bottomrule
\end{tabular}
\end{center}

\lstinline|Remap| is the quiet hero. Converting a health value to a bar width, a
mouse position to a world coordinate, a distance to a volume --- all one call:

\begin{lstlisting}
float volume = Remap(distance, 0.0f, 20.0f, 1.0f, 0.0f);   // near = loud
\end{lstlisting}

\subsection{Vector2, Vector3, Vector4}

The same set exists for each, with the obvious names:
\lstinline|Add|, \lstinline|Subtract|, \lstinline|Scale|, \lstinline|Multiply|,
\lstinline|Divide|, \lstinline|Length|, \lstinline|LengthSqr|,
\lstinline|Distance|, \lstinline|Normalize|, \lstinline|DotProduct|,
\lstinline|Lerp|, \lstinline|Negate|, \lstinline|Min|, \lstinline|Max|.

Plus a few that only make sense in one dimension:

\begin{itemize}[leftmargin=*,itemsep=2pt]
  \item 2D only --- \lstinline|Vector2Angle|, \lstinline|Vector2Rotate|
  \item 3D only --- \lstinline|Vector3CrossProduct|,
        \lstinline|Vector3RotateByQuaternion|
\end{itemize}

\keyidea{Use \lstinline|Vector3DistanceSqr| instead of
\lstinline|Vector3Distance| when you only want to compare two distances. Skipping
the square root is free, and comparing squared distances gives the same answer.
With hundreds of enemies it is a measurable win.}

\subsection{Matrix}

A \lstinline|Matrix| is a 4$\times$4 grid of floats that encodes a transform:
a translation, a rotation, a scale, or all three at once.

\begin{lstlisting}
Matrix t = MatrixTranslate(x, y, z);
Matrix r = MatrixRotateY(angle);
Matrix s = MatrixScale(2, 2, 2);

Matrix combined = MatrixMultiply(MatrixMultiply(s, r), t);   // scale, then rotate, then move
model.transform = combined;
\end{lstlisting}

\gotcha{Matrix multiplication is not commutative. Rotating then translating puts
the object somewhere completely different from translating then rotating. In
raylib's convention the transforms apply left to right, so read
\lstinline|MatrixMultiply(s, r)| as ``scale, then rotate''.}

You rarely build these by hand. Where they turn up is
\lstinline|model.transform|, which lets you position a model without touching the
draw call, and \lstinline|DrawMeshInstanced|, which takes an array of them.

\subsection{Quaternion}

A quaternion is four floats that represent a rotation. They exist because Euler
angles --- yaw, pitch, roll --- have two problems: gimbal lock (chapter 12), and
the fact that blending between two of them takes a scenic route.

\begin{lstlisting}
Quaternion q = QuaternionFromAxisAngle((Vector3){ 0, 1, 0 }, angle);
Vector3 spun = Vector3RotateByQuaternion(point, q);

Quaternion blended = QuaternionSlerp(from, to, t);    // smooth, shortest path
\end{lstlisting}

\lstinline|QuaternionSlerp| is the reason they exist: it interpolates between two
orientations along the shortest arc, at constant speed. That is what every
animation system and every smooth camera does internally.

For a first person camera, two angles are simpler and perfectly adequate --- keep
using yaw and pitch. Reach for quaternions when something needs to rotate freely
in all three axes, like a spaceship or a physics object.

\section{You do not have to use it}

raymath is a convenience, not a requirement. \lstinline|position.x += speed*dt|
is a perfectly good line of code, and for simple 2D work writing the arithmetic
out is often clearer than nesting three function calls.

Where it earns its place is 3D, where the vector operations get long and the
chance of a typo in the third component is high.

\tryit{Module 10 has two panels. In 2D, drag two vectors and watch add, subtract,
dot, angle, lerp and rotate update live. Press TAB for the 3D panel, where the
cross product and a quaternion rotation are drawn in
space.}{./cheatsheet/build.sh 10}
